\documentclass{article}
\usepackage[utf8]{inputenc}
\usepackage[T1]{fontenc}
\usepackage{geometry}
\usepackage{graphicx}
\usepackage{float}
\usepackage{hyperref}
\usepackage[english,indonesian]{babel}

\title{Optimasi Pemrosesan Template LaTeX Menggunakan Pendekatan Berbasis AI}
\author{Academic Support Agent}
\date{\today}

\begin{document}
\maketitle

\begin{abstract}

\end{abstract}

\chapter{PENDAHULUAN}


\chapter{TINJAUAN PUSTAKA}


% Original content (commented out):
% \input{config}

\begin{document}


\input{Chapters/1.Preliminary pages/1.Cover page}
\input{Chapters/1.Preliminary pages/2.Dedication}
\input{Chapters/1.Preliminary pages/3.Acknowledgements}
\input{Chapters/1.Preliminary pages/4.Abstract}
\input{Chapters/1.Preliminary pages/5.Glossary}
\input{Chapters/2.Unnumbered chapters/1.Introduction}
\input{Chapters/2.Unnumbered chapters/2.Problem statement}
\input{Chapters/3.Numbered chapters/1.Presentation of the corporate}
\input{Chapters/3.Numbered chapters/2.Project scope and used methodology}
\input{Chapters/3.Numbered chapters/3.Theoretical framework}
\input{Chapters/3.Numbered chapters/4.Operational framework}
\input{Chapters/3.Numbered chapters/5.Results}
\input{Chapters/3.Numbered chapters/6.Discussion}
\input{Chapters/4.Back matter/1.References}
\input{Chapters/4.Back matter/2.Appendices}

\end{document}


\end{document}